\begin{proof}

\begin{enumerate}[label={$\langle1\rangle$\arabic*.},leftmargin=*]
\item Let $\nu:T\to S$ denote this finite normalization morphism.
A continuous bijection from the compact $T$ to the Hausdorff
$S$ is a homeomorphism. It therefore induces an isomorphism
of fundamental groups and equality of Euler numbers. P4R.1
makes $\pi_1(T)$ virtually abelian.

\item The effective-conductor and canonical-torsion argument of
P4.12 applies to this same $T$ and proves $\kappa(T)=-\infty$.
The opened classification theorem used there leaves rational,
ruled, and class-VII minimal models, followed by point blowups.
Blowing up a smooth point preserves the fundamental group:
van Kampen applies to the common punctured four-ball with
simply connected boundary $S^3$. A ruled surface over a
genus-$g$ curve has the same fundamental group as that curve,
by the homotopy sequence of its $\mathbb P^1$-bundle.
For $g\ge2$ this group is not virtually abelian. Indeed a
finite-index subgroup corresponds to an unramified cover
of genus $1+d(g-1)>1$, and a genus-$h$ surface group maps
onto the nonabelian free group on $h$ generators by sending
the $b$ generators to the identity in its usual presentation.
Consequently $g\le1$.

\item A rational surface has positive Euler number. A ruled surface
over an elliptic curve, blown up $k$ times, has Euler number
$k$. A class-VII surface has Euler number $b_2\ge0$, and
each blowup increases it by one. These formulas give
$e(S)=e(T)\ge0$. Equality in the ruled case is exactly
$k=0$. Equality in class VII gives a minimal surface with
$b_2=0$. The conductor is nonempty because $S$ is singular
and nonnormal; P4.12 gives an effective conductor curve on
$T$. The precise class-VII classification and absence of
curves on Inoue surfaces verified in P4.8 therefore force
$T$ to be Hopf. The converse follows from the same Euler
formulas and the homeomorphism $\nu$.

\end{enumerate}
\end{proof}
